\chapter{Introduzione}
\section{Motivation and Index Construction}

While a vast literature has explored the effects of trade liberalization in developed economies, less is known about its consequences in developing countries, particularly concerning internal stability and conflict. In these contexts, trade represents a major channel through which global economic fluctuations are transmitted to local welfare and political stability \parencite{united2011review}. Existing research has often examined related mechanisms through the lens of commodity price shocks. For instance, \textcite{dube2013commodity} and \textcite{bruckner2010international} analyse how international price variations affect conflict outcomes, finding mixed results depending on the type of resource and the institutional environment. Other contributions have instead focused on the reverse relationship, investigating how conflict disrupts trade flows and economic integration \parencite{martin2008civil}. 
However, evidence on how trade shocks themselves, rather than price or conflict shocks—affect local stability remains limited, especially at the subnational level in low- and middle-income countries. Our analysis addresses this gap by examining how spatial exposure to trade integration shapes conflict dynamics within Nigeria, combining detailed information on port activity, global connectivity, and local socio-economic characteristics.
Our analysis builds on the intuition of \textcite{berman2015external}, who show that positive income shocks, such as those generated by trade, can reduce conflict by improving local economic opportunities. Conversely, trade can also exacerbate tensions when it increases the value of contestable rents or reinforces existing inequalities, as shown by \textcite{amodio2024trade}. These dual mechanisms, often referred to as the \textit{opportunity} or \textit{income} and \textit{rent-seeking} effects, suggest that the impact of trade on conflict is theoretically ambiguous and context-dependent.

Nigeria represents an ideal setting to study these dynamics. It is a coastal country with a diversified set of ports and abundant natural resources, where trade-related income shocks can have significant spatially heterogeneous effects. Its geography—marked by strong North–South disparities and localized ethnic divisions—makes it a relevant case for testing whether trade integration promotes peace or amplifies existing inequalities.

To empirically assess these mechanisms, we construct a spatial trade exposure index inspired by shift–share approaches. The index combines three main components:
\begin{enumerate}
	\item \textbf{Trade shocks}, measured as annual changes in port-level container throughput from IMF PortWatch;
	\item \textbf{Connectivity}, proxied by the Liner Shipping Connectivity Index (LSCI) from UNCTAD, which captures each port’s integration in global maritime networks;
	\item \textbf{Proximity}, expressed as the inverse distance between each ADM2 centroid and the nearest port (including Onne, added to complement the FAO–UNCTAD list).
\end{enumerate}

INDEX
This formulation captures how strongly each subnational unit is exposed to international trade shocks depending on its geographical proximity and the global relevance of its nearest ports. Regions located closer to active and well-connected ports are assumed to experience stronger transmission of trade-related income shocks.

Our contribution is twofold. First, we provide one of the first subnational evaluations of trade shocks and conflict in Africa, combining spatially detailed datasets on trade, conflict, infrastructure, and environmental factors. Second, we show that the effect of trade openness is highly heterogeneous: while it fosters stability in resource-constrained and less developed areas, it may reinforce instability in already prosperous or rent-dependent regions. These findings highlight the importance of considering spatial inequality in designing conflict-sensitive trade and infrastructure policies.
